\chapter{基于YOLO26与点云投影的吊装构件空间定位、预警与避障方法}\label{chap4_yolo26_projection_warning}

\section{基于YOLO26的吊装构件视觉感知模型与边缘端部署}\label{sec_ch4_yolo26_deploy}

\subsection{YOLO26模型原理与选型依据}\label{subsec_ch4_yolo26_selection}

本章首先将视觉侧任务重新界定为\textbf{吊装构件或吊物区域的实时识别与分割}，而不是将视觉网络直接用于人员或车辆的全流程轨迹跟踪。这样调整的原因在于，吊装作业的风险判定并不依赖图像平面中的目标外观变化本身，而依赖于吊装构件在三维空间中的位置、边界及其与障碍物之间的距离关系。视觉网络在整条链路中的作用，是为点云处理提供可靠的二维语义约束，使后续点云投影与三维定位能够围绕“吊装构件在哪里”这一核心问题展开，而不是在视觉侧提前承担全部运动理解任务。

在模型选择上，本文采用 YOLO26 分割模型作为吊装构件视觉感知的基础模型。选择 YOLO26 而非继续沿用旧版本检测模型，主要基于三个方面的考虑。其一，YOLO26 属于单阶段实时感知框架，前向推理链路较短，适合部署在 Jetson 边缘端以满足在线处理需求。其二，本研究更加关注\textbf{区域级语义结果}，即需要获得吊装构件在图像中的完整轮廓或掩码，而不仅是粗粒度的矩形框，因此分割输出天然比检测框更适合作为点云回查的空间约束。其三，从工程迁移成本看，现有系统已经完成了 YOLO26 的 ROS 封装、TensorRT 推理接入与分割结果发布，继续沿该路径完善论文论述，能够保持论文方法与代码实现之间的一致性。

需要强调的是，YOLO26 在本文中的定位不是孤立的识别器，而是整个雷视融合链路的语义起点。设输入图像为 $I_t$，YOLO26 分割网络输出为实例集合
\begin{equation}
\mathcal{S}_t=\{(M_i,c_i,s_i)\mid i=1,2,\dots,N_t\},
\end{equation}
其中 $M_i$ 表示第 $i$ 个实例的二值分割掩码，$c_i$ 表示类别标签，$s_i$ 表示置信度分数。对于本文当前阶段，类别集合以吊装构件或吊物相关类别为中心，具体的类别数、标注策略和训练细节仍需结合后续训练记录进一步补齐，因此正文中关于类别定义可保留为 \textbf{【待补充：类别集合与标注口径】}。这一处理能够保证方法框架先行明确，同时避免正文与真实训练记录不一致。

从系统角色划分来看，YOLO26 负责回答“哪些二维区域属于吊装构件”，LiDAR 与定位模块负责回答“这些区域在三维空间中对应哪些点、其空间中心和距离是多少”。这种分工符合吊装场景的工程特点：LiDAR 具备直接测距能力，视觉网络则更适合提供语义边界。二者的组合避免了仅依赖视觉深度估计所带来的尺度漂移问题，也避免了仅依赖点云聚类时难以稳定区分吊物与其他结构反射的局限。因此，本节的结论是，YOLO26 的价值不在于单独完成风险判断，而在于为后续空间定位构建稳定、可解释的二维语义入口。

\subsection{面向Jetson边缘端的TensorRT部署与FP16半精度量化}\label{subsec_ch4_tensorrt_fp16}

对于吊装安全预警系统而言，模型是否能够稳定运行在边缘端，往往比离线精度能否继续提升若干百分点更关键。原因在于吊钩端平台受体积、功耗和散热条件限制，感知算法必须在有限算力下维持可接受的吞吐率与时延波动范围。因此，本文在视觉侧采用“训练权重导出为 ONNX，再转换为 TensorRT engine，在 Jetson 平台上执行 FP16 半精度推理”的工程部署路径。这一路径已经在现有系统中形成清晰实现基础，也是本章中可以直接对应代码的部分。

从部署流程看，YOLO26 模型训练完成后，首先需要导出 ONNX 表达，以便与 TensorRT 推理后端对接；随后利用 TensorRT 对网络图进行层融合、内存规划和硬件相关优化，得到边缘端可直接加载的 engine 文件。当前工程的实现表明，YOLO26 ROS 节点在初始化阶段直接面向 TensorRT 推理后端进行封装，推理接口已经固定在分割任务场景下工作。与此同时，工程记录中还考虑了 engine 元数据兼容问题，即当模型文件在导出与加载链路之间存在额外头部或元信息时，需要在加载前做适配处理。该问题虽然属于实现细节，但在工程部署中非常典型，说明本文的系统并非停留在理论选型，而是已经进入可运行的边缘端落地阶段。

FP16 半精度量化的采用，主要是为了在不显著牺牲推理稳定性的前提下降低显存占用并提升吞吐率。设单帧推理耗时为 $t_{\mathrm{infer}}$，边缘端总处理周期为 $T$，则系统可持续在线运行的一个基本条件是
\begin{equation}
t_{\mathrm{infer}} + t_{\mathrm{pre}} + t_{\mathrm{post}} \leq T,
\end{equation}
其中 $t_{\mathrm{pre}}$ 与 $t_{\mathrm{post}}$ 分别表示图像预处理和后处理耗时。在 Jetson 平台中，若仍采用较高精度推理或不进行 TensorRT 图优化，则 $t_{\mathrm{infer}}$ 容易成为整条链路的主瓶颈。采用 FP16 后，视觉推理在整体时延预算中的占比得到压缩，从而为点云转换、投影、叠加绘制和结果发布留出更稳定的时延空间。当前阶段系统已经具备视觉链路的阶段耗时统计，因此第5章可以基于这些计时结果进一步展开实时性分析。

需要说明的是，本节可以明确写成“已有工程实现基础”，但不能夸大为“训练与部署全流程的每个细节都已在仓库中完整留存”。例如 engine 文件路径、训练平台细节、导出时的工作空间参数等，当前仓库并未全部保留，因此正文中宜采用“已完成部署链路验证，具体导出参数为 \textbf{【待补充：ONNX 导出与 trtexec 参数】}”的写法。这样既保留了系统已经落地到 Jetson 边缘端这一事实，也避免因细节记录不全而造成表述失真。

\begin{figure}[htbp]
	\centering
	\fbox{\parbox[c][5.2cm][c]{0.88\textwidth}{\centering
	占位图：YOLO26 模型训练、导出与 Jetson 边缘端部署流程图\\[0.5em]
	建议内容：工地实采数据整理与标注 $\rightarrow$ YOLO26 训练与验证 $\rightarrow$ ONNX 导出 $\rightarrow$ TensorRT engine 生成 $\rightarrow$ Jetson 端 ROS 节点加载与在线推理。}}
	\caption{YOLO26 模型训练、导出与边缘端部署流程占位图}
	\label{fig:ch4_yolo26_deploy_placeholder}
\end{figure}

\subsection{工地实采数据上的模型训练、验证与预测流程}\label{subsec_ch4_train_validate_predict}

从研究完整性看，视觉模型的构建不应只停留在“部署成功”这一结果，还应说明数据来源、训练组织方式以及在线预测流程。本研究已知的可靠事实是：YOLO26 模型所对应的数据来源于工地实际采集样本，模型当前在系统中承担吊装构件分割任务，并已经形成“输入图像、在线推理、输出分割结果、结果发布”的预测链路。相较于公开通用数据集，工地实采数据更能反映遮挡、背景杂乱、构件尺度变化以及照明波动等真实作业条件，因此其价值在于提高模型对目标场景的贴合程度。

在论文组织上，本小节宜分为“数据整理与标注、训练配置、验证评估、模型导出与部署验证”四个部分展开。当前仓库对最后一部分，即模型导出与部署验证，支撑较强；对前两部分，即类别定义、训练轮次、学习率、增强策略和训练平台，支撑尚不充分。因此正文应采取谨慎写法：先说明训练流程在研究方法中的必要位置，再把当前已经确认的事实写实，把尚未补齐的训练细节留作占位。例如，可写为“模型训练采用工地实采数据构建样本集，并完成面向吊装构件分割任务的训练与部署验证；具体类别规模、训练参数和验证指标记为 \textbf{【待补充：训练配置与精度结果】}”。

与训练部分相比，在线预测流程在当前系统中更适合直接展开。图像流进入 ROS 节点后，模型首先完成前向推理与实例分割；随后系统将分割掩码与原始图像尺寸对齐，并在需要时叠加投影结果与可视化信息，最终发布分割图像及其关联结果。这一流程说明，本文的视觉侧并不是孤立输出一个离线精度指标，而是以在线消息流的形式进入后续空间定位链路。因此，在论文写法上，预测流程应被明确界定为“从视觉感知到空间定位的前端入口”，而训练章节则作为该入口能够成立的模型基础。这样的安排既符合实际代码结构，也使第4章后续各节之间的逻辑更加紧密。

\section{点云与图像关联方法及吊装构件空间定位}\label{sec_ch4_projection_localization}

\subsection{图像-点云时空同步与坐标系统一}\label{subsec_ch4_sync_frames}

视觉分割结果只有在与点云保持时间对齐和坐标一致的前提下，才能转化为可信的三维定位结果。因此，本节首先处理多传感器融合中的基础问题，即图像帧与点云帧如何同步，以及各类数据如何统一到一致的参考坐标系。当前系统中，默认输入已经采用 leveled 图像和 leveled 点云，这意味着前端传感器数据在进入投影链路之前，已经经过姿态或参考系上的预处理，从而降低了吊钩端摆动对投影稳定性的直接影响。

在时间同步方面，现有实现采用近似时间同步策略，而不是强求两类传感器严格同一时刻采样。这一选择的工程依据在于，相机、LiDAR、ROS 消息传输和模型推理本身都存在时延与抖动；若只接受严格相同时间戳的数据对，实际系统中会频繁发生匹配失败。近似时间同步通过为图像帧和点云帧设置一个可接受的最大时间间隔，使系统能够在小范围时间误差内继续完成投影与关联。设图像时间戳为 $t_{\mathrm{img}}$，点云时间戳为 $t_{\mathrm{pc}}$，同步约束可写为
\begin{equation}
|t_{\mathrm{img}}-t_{\mathrm{pc}}| \leq \Delta t_{\max},
\end{equation}
其中 $\Delta t_{\max}$ 表示允许的近似同步窗口。该参数既影响投影稳定性，也影响系统可用性，第5章可将其作为时延与效果分析的控制变量之一。

在坐标系统一方面，当前实现采用“TF 优先、静态配置兜底”的方式构建投影模型。也就是说，当运行时 TF 能够给出相机与点云坐标系之间的实时变换时，系统优先使用 TF 结果；若 TF 不可用或暂时查询失败，则回退到 YAML 中预先标定的静态外参。这种设计的意义在于，一方面保证了在线运行的灵活性，另一方面避免因单一外参来源失效而导致整个投影链路中断。对论文表述而言，这种机制应写成“增强系统稳健性的工程处理”，而非单纯的实现细节，因为它直接关系到实际作业环境下投影结果的连续性。

综上，时空同步与坐标统一构成了第4.2节全部后续方法的前提。如果这一基础链路不稳定，那么即使视觉分割或点云质量本身较好，最终得到的投影对应关系也会出现错位。因此，本小节的结论是：吊装构件空间定位不是从投影公式开始，而是从同步与坐标语义的统一开始；只有在此基础上，视觉语义与点云几何之间的对应关系才具有工程可信性。

\subsection{基于分割掩码的3D到2D投影与点云集合回查}\label{subsec_ch4_mask_projection}

当前系统的核心实现并不是“由单个图像像素反求唯一三维点”，而是采用更符合 LiDAR 几何特性的路径，即先将三维点云逐点投影到二维图像平面，再利用分割掩码筛选属于吊装构件实例的点云集合。该方法的关键优势在于，LiDAR 已经直接给出了三维点的空间坐标，视觉网络只需提供二维语义边界，二者通过投影关系建立联系即可，无需额外构造不稳定的单目深度反演过程。

设 LiDAR 坐标系下某一点为 $\mathbf{p}_L=[x_L,y_L,z_L,1]^\mathsf{T}$，相机内外参分别为 $\mathbf{K}$ 与 $\mathbf{T}_{CL}$，则其投影到图像平面的过程可表示为
\begin{equation}
\mathbf{p}_C=\mathbf{T}_{CL}\mathbf{p}_L,
\qquad
\begin{bmatrix}
u'\\v'\\w'
\end{bmatrix}
=\mathbf{K}
\begin{bmatrix}
x_C\\y_C\\z_C
\end{bmatrix},
\qquad
u=\frac{u'}{w'},\; v=\frac{v'}{w'}.
\end{equation}
对所有满足视场和深度约束的 LiDAR 点执行上述投影后，可得到它们在图像中的像素位置 $(u,v)$。若某点投影落入第 $i$ 个分割掩码 $M_i$ 对应的前景区域，则将该点归入实例点集 $\mathcal{P}_i$，即
\begin{equation}
\mathcal{P}_i=\{\mathbf{p}_L\mid (u,v)\in M_i\}.
\end{equation}
这样，二维语义分割结果就被转换为一个个具有明确三维坐标的点云子集，为后续空间定位提供了直接输入。

从实现角度看，这一路径已经在现有代码中形成清晰结构：投影模型负责封装内参与外参，投影函数负责完成三维点到二维像素的映射，而绘制与实例处理函数则在遍历分割对象时，根据掩码回查点云并组织可视化输出。论文在此处应强调，这种方法不是对每个像素进行“求逆”，而是对每个已有三维点进行“求正向投影并判定其是否属于目标区域”。这一点非常重要，因为它决定了方法的误差来源主要是外参误差、同步误差与掩码边界误差，而不是深度反演误差。

此外，该方法也非常符合吊装场景的任务目标。对于安全预警而言，系统关心的是“吊物相关点集在三维空间中的位置与边界”，而不是图像中的纹理细节本身。分割掩码能够比矩形框更精确地约束目标区域，从而减少背景点被错误回收到目标点集中的概率。因此，本小节可以作为整章中最关键的已实现方法之一：它直接把视觉语义和点云几何连接起来，是后续空间定位、距离计算和可视化输出成立的技术基础。

\begin{figure}[htbp]
	\centering
	\fbox{\parbox[c][5.8cm][c]{0.9\textwidth}{\centering
	占位图：基于分割掩码的 3D 点投影到 2D 图像并回查目标点云流程图\\[0.5em]
	建议内容：输入图像与点云 $\rightarrow$ 相机内外参投影 $\rightarrow$ 3D 点落入图像平面 $\rightarrow$ mask 判定前景像素 $\rightarrow$ 回查实例点集 $\rightarrow$ 生成彩色点云与质心。}}
	\caption{基于分割掩码的点云投影回查与实例点集提取占位图}
	\label{fig:ch4_projection_workflow_placeholder}
\end{figure}

\subsection{吊装构件三维位置估计与距离基准构建}\label{subsec_ch4_3d_localization}

当目标实例对应的点云子集 $\mathcal{P}_i$ 被提取出来之后，下一步任务不再是判断“是否识别到目标”，而是从这些点中构建能够服务于预警计算的三维空间表达。当前系统已经具备基于实例点集计算三维质心的实现基础，因此可以将质心作为最直接的吊装构件位置估计量。设实例点集包含 $n_i$ 个三维点，则其质心可定义为
\begin{equation}
\bar{\mathbf{p}}_i=\frac{1}{n_i}\sum_{k=1}^{n_i}\mathbf{p}_{i,k}.
\end{equation}
该结果能够反映吊装构件在当前时刻的整体空间位置，并可与图像投影结果一同用于可视化展示与后续模块对接。

然而，仅使用质心并不足以完整描述预警所需的空间关系。对于吊装作业而言，真正参与风险判定的并不是单一点位置，而是吊装构件与周边障碍物之间的最小安全边界。因此，论文在这里需要进一步提出“距离基准构建”的概念，即在三维点集基础上构建更适合安全判定的空间表达方式。可选方案至少包括三类。第一类是\textbf{质心表达}，其优点是简单、稳定，适合做目标跟踪与全局位置描述；缺点是当目标尺寸较大或形状细长时，质心无法直接反映最近碰撞边界。第二类是\textbf{空间包围范围表达}，例如最小包围盒或轴对齐边界盒，可更好地描述目标几何尺寸；第三类是\textbf{最近邻边界点表达}，直接利用目标点集与障碍物点集之间的最近距离作为风险计算输入，其保守性更强，但对噪声也更敏感。

结合当前代码现状，本文可以将“质心已实现，空间包络与最近边界距离为后续补充设计”写清楚。也就是说，现有系统已经完成了目标实例点云的提取与质心计算，能够支撑“吊装构件三维位置估计”这一论断；但尚未形成专门的吊物空间包络发布接口，因此更细粒度的边界建模仍应按论文方法设计进行展开。这样写的优点在于，既不弱化当前系统已经具备的三维定位基础，又不把尚未落地的空间边界建模误写成已完成模块。

基于这一处理，本小节的结论可以表述为：当前系统已经从二维分割结果成功过渡到三维目标点集与质心位置，证明了吊装构件空间定位链路的可行性；进一步面向预警与避障的距离基准构建，则将在后续章节中通过分层安全边界和动态风险评估继续完善。

\subsection{雷视融合数据接口设计与下游输出组织}\label{subsec_ch4_data_interface}

为了使视觉分割与点云投影结果真正服务于预警模块，仅有图像叠加效果图还不够，还需要把结果组织为可供下游调用的结构化数据接口。当前系统已经具备两类直接输出：一类是包含分割与投影效果的图像结果，用于人机可视化与离线分析；另一类是语义着色点云，用于展示哪些三维点被识别为吊装构件相关点。对于论文写作而言，这两类输出说明第4.2节并非停留在理论公式，而已经形成了可验证的工程结果。

进一步面向预警链路，本文建议将输出接口组织为至少三层。第一层是\textbf{可视化层}，包括投影叠加图像与语义着色点云，其作用是帮助观察系统运行状态与定位效果。第二层是\textbf{目标描述层}，包括目标时间戳、类别标签、置信度、三维质心和实例点数等字段，用于把感知结果转化为程序可消费的对象状态。第三层是\textbf{风险计算准备层}，在目标描述基础上继续补充最小距离、空间包络、速度估计来源和数据有效性标志位，为后续预警模块提供统一输入。

由于当前仓库中尚未定义完整的统一消息类型，因此本小节应明确写成“接口设计与输出组织方案”。建议输出结构至少包含
\begin{equation}
\mathbf{z}_i=\{t_i,c_i,s_i,\bar{\mathbf{p}}_i,d_i,\eta_i\},
\end{equation}
其中 $t_i$ 为时间戳，$c_i$ 为类别标签，$s_i$ 为置信度，$\bar{\mathbf{p}}_i$ 为三维位置估计，$d_i$ 为最小距离或边界距离，$\eta_i$ 为数据有效性标志。若后续需要承接动态风险评估，还可以在接口中补充速度、预测轨迹句柄或包络信息。通过这一设计，感知模块输出将不再只是“好看的结果图”，而是能够真正进入风险计算与控制决策的规范化对象状态。

\section{面向吊装作业的多层级预警链路设计}\label{sec_ch4_warning_design}

\subsection{基于空间距离的静态风险阈值卡控}\label{subsec_ch4_static_distance_warning}

在完整的吊装安全预警链路中，最基础的一层判断应当来自空间距离阈值卡控。其核心思想是：不论是否存在复杂动态目标，只要吊装构件与周边静态障碍物之间的最小空间距离低于安全边界，就应触发相应等级的风险提示。与传统仅基于图像框位置或主观经验判断不同，本研究强调以第4.2节得到的三维空间定位结果为基础建立安全边界，从而使预警判定建立在几何量之上。

设吊装构件状态为 $\mathcal{O}$，静态障碍物集合为 $\mathcal{E}=\{E_1,E_2,\dots,E_m\}$，则吊装构件与环境之间的最小距离可表示为
\begin{equation}
d_{\mathrm{static}}=\min_{E_j\in\mathcal{E}}\mathrm{dist}(\mathcal{O},E_j).
\end{equation}
在此基础上，可将静态预警划分为提示区、警戒区和危险区三个层级，分别对应阈值 $D_1>D_2>D_3$。当 $d_{\mathrm{static}}<D_1$ 时，系统进入一级提示；当 $d_{\mathrm{static}}<D_2$ 时，进入二级报警；当 $d_{\mathrm{static}}<D_3$ 时，进入三级危险状态。该设计的意义在于，它为后续更复杂的动态升级规则提供了一个几何上的基础层，使系统能够先快速筛除明显安全的场景。

需要指出的是，当前仓库中旧版基于网格风险区的逻辑已经被移除，新的分层距离卡控尚未形成完整实现。因此，本小节必须明确写成\textbf{论文方法设计}，而不是既有代码功能描述。这样处理并不削弱章节逻辑，反而更符合论文的真实贡献边界：已有代码提供了吊装构件三维位置与彩色点云基础，而本文在此基础上系统化提出了可直接工程实现的静态风险分层判定方法。对应的阈值数值、标定方式和实验结果将在第5章通过 \textbf{【待补充：距离阈值标定与误报统计】} 进一步完善。

\subsection{动态点检测、聚类与障碍物状态表达}\label{subsec_ch4_dynamic_candidate}

静态距离阈值只能描述环境边界对吊装构件的约束，而无法充分反映人员、车辆或其他运动体带来的时间维度风险。因此，预警链路还需要承接动态点检测模块输出的障碍物候选状态。当前系统已经形成较完整的上游动态链路，包括 LiDAR 数据输入、定位建图、动态点识别、欧式聚类和目标中心提取。这意味着，论文在本小节中可以把“动态障碍物候选生成器”写成已具备实现基础的模块，而不是纯粹的构想。

从功能上看，动态点检测模块的目标不是直接输出最终风险等级，而是把杂散的动态点组织成目标级状态量。设经过动态筛选后的点集为 $\mathcal{Q}_t$，经聚类后得到候选目标集合
\begin{equation}
\mathcal{C}_t=\{C_1,C_2,\dots,C_{n_t}\},
\end{equation}
其中每个聚类对象 $C_k$ 至少应包含几何中心、点数规模、空间范围以及时间戳等属性。只有当动态点从“离散点集合”上升为“对象状态”后，后续的速度估计、轨迹预测和风险评估才具备统一输入接口。

在论文表达中，本小节还应明确区分“候选生成”和“风险确认”这两个阶段。动态点检测模块只说明某处出现了具有运动特征的目标，并不意味着该目标一定构成与吊装构件有关的风险。真正的风险判断还需要结合吊装构件位置、相对运动方向、未来交汇趋势以及安全距离阈值综合计算。因此，动态点检测在全系统中的合理定位应是“为预警模块提供动态障碍物候选状态”，而不是直接承担预警决策职责。这样的表述既符合当前代码组织，也有利于后续在第4.3.4节引入更高层的风险升级规则。

\subsection{基于卡尔曼滤波的未来运动趋势预测}\label{subsec_ch4_kalman_prediction}

在获得动态障碍物候选目标之后，系统需要进一步估计其未来短时运动趋势，以判断其是否会与吊装构件形成持续接近关系。当前代码已经实现二维常速度模型的卡尔曼滤波器，因此本小节可以直接建立在现有算法基础之上展开。设动态目标在平面内的状态向量为
\begin{equation}
\mathbf{x}_k=[p_x,p_y,v_x,v_y]^\mathsf{T},
\end{equation}
则在采样时间间隔 $\Delta t$ 下，常速度模型的状态转移可写为
\begin{equation}
\mathbf{x}_{k+1}=\mathbf{F}(\Delta t)\mathbf{x}_k+\mathbf{w}_k,
\qquad
\mathbf{F}(\Delta t)=
\begin{bmatrix}
1&0&\Delta t&0\\
0&1&0&\Delta t\\
0&0&1&0\\
0&0&0&1
\end{bmatrix},
\end{equation}
其中 $\mathbf{w}_k$ 为过程噪声。观测量通常由动态聚类中心给出，对应位置更新项 $[p_x,p_y]^\mathsf{T}$。

该模型的优点在于形式简单、参数可控，能够在动态目标短时运动较平滑的前提下稳定输出位置和速度联合估计。当前系统默认预测时域为 2.0 s，预测步长为 0.2 s，并能够发布未来轨迹可视化结果。对论文而言，这些实现细节具有直接价值，因为它们表明系统并非只在当前时刻做几何检测，而是已经具备面向短时未来的趋势估计能力。第5章中的动态场景实验也可以围绕这一预测时域和步长展开分析。

需要注意的是，本小节的目标不是把卡尔曼滤波本身写成主要创新，而是说明它在整条预警链路中的功能位置。卡尔曼预测输出的未来轨迹并不直接等于碰撞风险结论，而是作为第4.3.4节相向运动升级判据和最近接距离估计的前置输入。换言之，预测模块回答的是“目标未来可能怎样运动”，而风险评估模块回答的是“这种运动是否会演化为危险接近”。把二者区分开来，有助于保持论文章节逻辑清晰，也能准确反映当前代码已经完成的部分与论文进一步补充的部分之间的关系。

\subsection{面向相向运动场景的风险升级判据与三级预警状态机}\label{subsec_ch4_risk_upgrade}

在静态距离阈值卡控与动态趋势预测的基础上，本文进一步提出面向相向运动场景的风险升级方法。其基本出发点是：仅依据当前时刻的瞬时距离，无法充分反映“虽然距离尚可，但正在快速接近”的危险状态；对于吊钩或吊物与动态障碍物相向运动的场景，若仍等到距离阈值被击穿后再触发高等级报警，则预警提前量可能不足。因此，需要把相对速度方向、未来最近接距离以及预测交汇趋势引入分级逻辑。

设吊装构件与某一动态目标的相对位置和相对速度分别为 $\mathbf{r}(t)$ 与 $\mathbf{v}_{\mathrm{rel}}(t)$，则可定义接近判据
\begin{equation}
\gamma(t)=\mathbf{r}(t)^\mathsf{T}\mathbf{v}_{\mathrm{rel}}(t).
\end{equation}
当 $\gamma(t)<0$ 时，表示二者处于相向接近或整体收敛趋势；当 $\gamma(t)\geq 0$ 时，表示相对距离不再缩小。在此基础上，还可引入未来预测窗口内的最近接距离
\begin{equation}
d_{\min}^{\mathrm{pred}}=\min_{\tau\in[0,T_p]}\|\mathbf{r}(t+\tau)\|,
\end{equation}
以及时间到最近接近或 TTC 指标。若同时满足“当前距离已进入警戒区”“相对运动方向为收敛”“预测最近接距离持续减小”三项条件，则系统应当把预警等级从静态层级进一步上调。

基于上述判据，本文建议构建三级预警状态机：一级提示主要用于提示操作员关注接近趋势；二级报警用于表明需要降低速度或限制部分运动方向；三级紧急用于表示已出现高碰撞风险，应触发停止或紧急制动。状态转移可由静态距离等级 $L_s$ 与动态趋势等级 $L_d$ 综合决定，即
\begin{equation}
L=\max(L_s,L_d).
\end{equation}
其中 $L_s$ 由静态距离阈值给出，$L_d$ 由相对速度方向、TTC、预测最近接距离等共同给出。为抑制状态频繁跳变，状态机还应设置进入与退出滞回、最小保持时间以及异常观测屏蔽规则。

需要特别说明的是，当前仓库中尚未完成 TTC 计算、相向运动升级规则和完整三级状态机实现，因此本小节属于\textbf{论文提出的系统化设计方案}。之所以仍需在第4章中详细展开，是因为现有动态检测与卡尔曼预测已经提供了足够坚实的前置基础，使这些设计并非脱离实现条件的空想，而是能够自然承接当前系统并向完整预警闭环扩展的下一步工作。

\begin{figure}[htbp]
	\centering
	\fbox{\parbox[c][5.4cm][c]{0.88\textwidth}{\centering
	占位图：多层级预警链路与三级状态机示意图\\[0.5em]
	建议内容：静态距离阈值卡控、动态候选生成、卡尔曼预测、相向运动升级判据、一级提示、二级报警、三级紧急及状态转移条件。}}
	\caption{多层级预警链路与三级预警状态机占位图}
	\label{fig:ch4_warning_state_machine_placeholder}
\end{figure}

\subsection{预警信息发布机制与系统接口组织}\label{subsec_ch4_warning_interface}

当预警等级被确定之后，系统需要把结果以统一、解耦的方式发布给上层人机界面或控制模块。因此，预警信息发布机制本身也是第4章必须明确的一部分。若没有统一接口，即使感知和判级方法已经成立，系统仍难以实现跨模块协同和后续控制对接。

本文建议将预警输出组织为两类消息。第一类是\textbf{风险状态消息}，主要服务于显示、日志记录和系统监控，字段可包括风险等级、触发对象 ID、对象类别、当前最小距离、预测最近接距离、相对速度、有效时间窗和数据有效性标志。第二类是\textbf{控制建议消息}，主要服务于第4.4节的避障执行模块，字段可包括建议动作类型、限速比例、限向标志位、急停标志位以及建议保持时间等。前者强调“系统判断到了什么风险”，后者强调“系统建议采取什么动作”，二者在接口上应当解耦，以便不同权限层级的系统按需订阅。

形式化地，可将风险状态消息记为
\begin{equation}
\mathbf{w}=\{L,\mathrm{id},c,d_{\mathrm{cur}},d_{\min}^{\mathrm{pred}},v_{\mathrm{rel}},T_p,\eta\},
\end{equation}
其中 $L$ 为风险等级，$\eta$ 为有效性标志。若后续系统接入控制接口，则再基于 $\mathbf{w}$ 生成控制建议消息 $\mathbf{u}_{\mathrm{adv}}$。这样设计的好处在于，即使控制模块尚未真正接入，预警链路本身也能独立运行和验证；而一旦需要向控制闭环扩展，也无需重构前端感知与判级输出格式。

当前仓库尚未定义统一预警消息类型，因此本小节同样应按设计方案写入论文。但与完全空白的设计不同，这里的接口定义是建立在已有图像、点云和动态轨迹等中间结果发布机制之上的，只是缺少最后一步标准化组织。因此，本小节的作用，是把现有“多个中间结果话题”进一步提升为“统一预警接口”的工程设计，为第4.4节控制接口的引入搭建数据基础。

\section{预警后的避障策略与控制接口设计}\label{sec_ch4_avoidance_design}

\subsection{从感知到控制的闭环避障架构}\label{subsec_ch4_closed_loop_arch}

为了使整套系统不止停留在“看见风险并发出提示”，还能够向安全辅助控制扩展，论文需要进一步说明从感知到控制的闭环架构。本文建议将系统组织为四层：第一层为视觉识别与吊物空间定位层，负责输出吊装构件状态；第二层为环境感知与动态候选层，负责输出静态障碍和动态目标状态；第三层为风险评估与预警分级层，负责融合距离、速度和趋势信息形成风险等级；第四层为控制接口与动作执行层，负责把风险等级转化为具体动作建议或控制指令。

这一架构的关键意义在于把“预警”和“避障”区分为两个不同层级的任务。预警层回答的是“当前风险有多高”，控制层回答的是“针对该风险应执行什么动作”。前者更强调感知可信性与判级逻辑，后者更强调接口规范、动作权限和执行安全性。将二者解耦能够保证系统在尚未接入实机控制的阶段，先完成预警链路验证；同时也为后续接入真实塔吊控制系统预留了清晰边界。

需要明确的是，当前仓库已经较为完整地覆盖了感知基础和动态趋势预测，但尚未形成统一控制接口与动作执行闭环。因此，本小节应采用“架构设计”写法，而不是陈述已完成控制模块。这样既不会夸大系统现状，也能体现论文在系统完整性层面的设计贡献，即把感知、预警和控制放在同一架构中讨论，而不是把控制问题完全留到论文之外。

\subsection{分级避障动作策略与控制接口约定}\label{subsec_ch4_control_strategy}

在控制接口层，动作策略应与预警等级一一对应，而不是由控制模块重新解释上游风险语义。为此，本文建议采用分级动作约定。一级提示对应观察与人工注意，不主动干预执行机构；二级报警对应限制部分运动方向或降低速度，例如降低起升速度、限制向障碍物方向的回转或行走；三级紧急对应硬停或软停指令，用于在高碰撞风险场景下优先保护人员与设备安全。这种分级动作机制能够保证风险等级与控制强度之间存在清晰映射关系。

若把控制建议表示为 $\mathbf{u}_{\mathrm{adv}}$，则其字段可组织为
\begin{equation}
\mathbf{u}_{\mathrm{adv}}=\{L,a,\rho,\delta,e,\tau_h\},
\end{equation}
其中 $L$ 为风险等级，$a$ 为动作类型，$\rho$ 为限速比例，$\delta$ 为受限方向或自由度标志，$e$ 为急停标志位，$\tau_h$ 为建议保持时间。通过这一设计，上层控制器可以根据自身权限和安全策略决定是否直接执行该建议，或仅将其作为辅助决策输入。也就是说，预警系统提供的是“可执行建议”，而不是强行越权接管控制。

对于相向接近和预测最近接距离持续减小的场景，应提高动作约束级别。例如，当系统判定动态目标与吊装构件存在持续收敛趋势，即使当前距离尚未进入最危险阈值，也可提前触发二级限速；当三级危险状态已成立时，则不再只给出提示，而应明确发布停止相关自由度的命令建议。该处理体现了本文设计中“动态趋势优先于单时刻距离”的安全思想。

当前仓库中尚未发现完整的控制接口实现，因此本小节必须明确写成\textbf{接口设计与动作约定}。但这一设计并非凭空提出，而是直接承接第4.3节所定义的预警输出字段，因此其工程实现路径是清晰的：只要上层控制系统能够接收标准化的风险状态和动作建议消息，就能进一步把论文中的设计转化为原型级闭环系统。

\begin{figure}[htbp]
	\centering
	\fbox{\parbox[c][5.6cm][c]{0.9\textwidth}{\centering
	占位图：从感知到控制的闭环避障架构图\\[0.5em]
	建议内容：视觉分割与空间定位层、静态/动态环境感知层、风险评估与预警分级层、控制建议输出层，以及与操作员或控制器之间的接口关系。}}
	\caption{从感知到控制的闭环避障架构占位图}
	\label{fig:ch4_closed_loop_arch_placeholder}
\end{figure}

\subsection{避障策略伪代码设计}\label{subsec_ch4_avoidance_pseudocode}

为突出整套预警与避障链路的可实现性，本文将核心执行逻辑进一步概括为伪代码式流程。其重点不在于展示复杂程序结构，而在于清楚说明“距离阈值卡控优先、动态趋势升级、控制动作分级执行”的主逻辑。按照第4章前述设计，系统每个处理周期都应先获取吊装构件状态、静态障碍物集合和动态轨迹预测结果，再统一完成风险判定与动作生成。

\begin{quote}
\ttfamily
Input:\\
suspended\_object\_state: 吊装构件三维位置、边界与速度信息\\
static\_obstacles: 静态障碍物集合\\
dynamic\_tracks: 动态目标当前状态与预测轨迹\\
threshold\_config: 距离阈值、TTC阈值与动作参数\\[0.3em]
Process:\\
1. 计算吊装构件与静态障碍物的最小距离 d\_static\_min；\\
2. 计算吊装构件与动态目标的相对运动关系，得到最近接距离、相对速度与 TTC；\\
3. 根据静态距离阈值得到等级 L\_s；\\
4. 根据动态趋势升级规则得到等级 L\_d；\\
5. 取 warning\_level = max(L\_s, L\_d)；\\
6. 若 warning\_level 为一级，则发布提示消息；\\
7. 若 warning\_level 为二级，则发布报警消息并生成限速或限向建议；\\
8. 若 warning\_level 为三级，则发布紧急消息并生成停止建议。\\[0.3em]
Output:\\
warning\_message, control\_advice
\end{quote}

上述流程说明，整套系统并不是由某一个复杂模型一次性完成所有判断，而是由“感知定位、风险计算、动作生成”三个相对独立但接口清晰的阶段构成。其优点在于结构可解释、便于调试，也更适合在工程上逐步迭代。例如，在当前阶段，系统已经能够稳定支撑前两个阶段中的大部分感知与定位能力；随着后续预警状态机和控制接口逐步补全，该流程即可自然扩展为更完整的安全辅助闭环。

本节的结论是，第4章从 YOLO26 视觉分割出发，经由点云投影回查建立吊装构件的三维空间定位基础，再进一步延伸到多层级预警和控制接口设计，构成了一条由已实现基础和论文方案设计共同支撑的完整技术路线。这种写法既忠实于当前代码实际，又能够体现论文对系统闭环性的提升和方法框架的完整性。